% !TeX program = pdflatex
% Changhao Li -- Research-focused industry CV
% Compile twice with pdfLaTeX; no .cls, .bib, images, or external font files needed.
% Research narrative draft revised on 11 October 2026.
% Inline draft comments identify personal contributions awaiting confirmation.
\RequirePackage{fix-cm}
\documentclass[11pt,letterpaper]{article}

% ======================= STYLE / TEMPLATE ================================
\usepackage[top=0.58in,bottom=0.58in,left=0.66in,right=0.66in,
            headheight=12pt,headsep=10pt,footskip=23pt]{geometry}
\usepackage[T1]{fontenc}
\usepackage[utf8]{inputenc}
\usepackage{tgpagella}
\usepackage[scale=0.95]{tgheros}
\usepackage{iftex}
\ifPDFTeX
  \input{glyphtounicode}
  \pdfgentounicode=1
\fi
\usepackage[protrusion=true,expansion=false]{microtype}
\usepackage{ragged2e}
\usepackage{xcolor}
\definecolor{Accent}{HTML}{234E64}
\definecolor{Body}{HTML}{20252B}
\definecolor{Muted}{HTML}{4D5963}
\definecolor{Rule}{HTML}{BBCBD3}
\usepackage{enumitem}
\usepackage{needspace}
\usepackage{fancyhdr}
\usepackage{lastpage}
\usepackage{hyperref}
\hypersetup{unicode=true,colorlinks=true,urlcolor=Accent,linkcolor=Accent,
  pdftitle={Changhao Li - Curriculum Vitae},pdfauthor={Changhao Li},
  pdfsubject={Agentic recursive self-improvement, AutoResearch, and long-horizon problem solving},
  pdfkeywords={Changhao Li, Georgia Tech, LLM, reinforcement learning, agents},
  pdfdisplaydoctitle=true,pdfborder={0 0 0}}
\urlstyle{same}
\setlength{\parindent}{0pt}
\setlength{\parskip}{0pt}
\setlength{\emergencystretch}{1.5em}
\hyphenpenalty=10000
\widowpenalty=10000
\clubpenalty=10000
\raggedbottom
\AtBeginDocument{\fontsize{11}{13.4}\selectfont\color{Body}\RaggedRight\hyphenpenalty=10000}

% Change these fields when reusing the template.
\newcommand{\CVName}{Changhao Li}
\newcommand{\CVUpdated}{October 2026}
\newcommand{\self}{\textbf{Changhao Li}}
\newcommand{\equal}{\textsuperscript{*}}
\newcommand{\second}{\textsuperscript{\textdagger}}
\newcommand{\sep}{\enspace\textcolor{Muted}{\textbar}\enspace}
\newcommand{\contactsep}{\enspace\textcolor{Muted}{\textperiodcentered}\enspace}
\newcommand{\smalltext}[1]{{\fontsize{10.2}{12.4}\selectfont #1\par}}
\newcommand{\meta}[1]{{\sffamily\fontsize{9.4}{11.6}\selectfont\color{Muted}#1\par}}
\newcounter{cvsec}
% Larger uppercase headings distinguish sections from body text and topic labels.
\newcommand{\cvsection}[1]{%
  \par\addvspace{12pt}\Needspace{4\baselineskip}\stepcounter{cvsec}%
  \pdfbookmark[1]{#1}{section\arabic{cvsec}}%
  {\sffamily\bfseries\fontsize{13}{15.5}\selectfont\color{Accent}\MakeUppercase{#1}\par}%
  \nobreak\vspace{3pt}{\color{Rule}\hrule height 0.55pt}\nobreak\vspace{5pt}}
% These are text-flow rows, not tables or separate columns.
\newcommand{\cvrow}[2]{%
  \par\noindent\textbf{#1}\hfill{\sffamily\fontsize{9.7}{11.7}\selectfont\color{Muted}#2}\par}
\newcommand{\job}[3]{%
  \Needspace{3\baselineskip}%
  \noindent\textbf{#1}\enspace---\enspace #2\hfill{\sffamily\fontsize{9.7}{11.7}\selectfont\color{Muted}#3}\par}
\newlength{\projectyearwidth}
\setlength{\projectyearwidth}{3em}
% A muted year gutter; wrapped titles align with the project name.
\newcommand{\project}[3]{%
  \par\addvspace{7pt}\Needspace{5\baselineskip}%
  \begingroup\hangindent=\projectyearwidth\hangafter=1
  \noindent\makebox[\projectyearwidth][l]{\sffamily\fontsize{9.7}{11.7}\selectfont\color{Muted}#1}%
  \textbf{#2}\enspace{\fontsize{10.2}{12.4}\selectfont #3}\par
  \endgroup\nobreak\vspace{1.5pt}}
\newenvironment{cvitems}{%
  \begin{itemize}[leftmargin=1.1em,label=\textbullet,labelsep=0.5em,
    topsep=1pt,itemsep=2.2pt,parsep=0pt,partopsep=0pt]\raggedright}{\end{itemize}}
\newenvironment{projectitems}{%
  \begin{itemize}[leftmargin=\dimexpr\projectyearwidth+1.1em\relax,label=\textbullet,labelsep=0.5em,
    topsep=1pt,itemsep=2.2pt,parsep=0pt,partopsep=0pt]\raggedright}{\end{itemize}}
% Publication: full title, full author list, venue/status, paper URL, code URL.
% Empty code URL is allowed. No per-publication abstract is printed.
% Topic labels are subordinate to the main section heading.
\newcommand{\pubtheme}[1]{%
  \par\addvspace{10pt}\Needspace{7\baselineskip}%
  \noindent{\bfseries\itshape\color{Accent}#1}\enspace
  {\color{Rule}\leaders\hrule height 0.35pt\hfill\kern0pt}\par
  \nobreak\vspace{1pt}}
\newcounter{pub}
\newcommand{\publication}[5]{%
  \par\addvspace{7.5pt}\Needspace{4\baselineskip}\stepcounter{pub}%
  \begingroup\fontsize{11}{13.3}\selectfont\raggedright
  \hangindent=1.9em\hangafter=0
  \noindent\llap{\makebox[1.9em][l]{\sffamily\fontsize{9.1}{11}\selectfont\color{Muted}[\arabic{pub}]}}%
  {\textbf{#1}}\par\nobreak
  \hangindent=1.9em\hangafter=0\noindent #2\par\nobreak
  \hangindent=1.9em\hangafter=0\noindent
  {\sffamily\fontsize{9.8}{11.8}\selectfont\color{Accent}\textbf{#3}\enspace\href{#4}{[Paper]}%
  \if\relax\detokenize{#5}\relax\else\enspace\href{#5}{[Code]}\fi}\par
  \endgroup}
\newcommand{\software}[3]{%
  \par\addvspace{3pt}\Needspace{3\baselineskip}%
  \smalltext{\textbf{\href{#2}{#1}}\enspace---\enspace #3}}
\newcommand{\award}[2]{%
  \smalltext{#1\hfill{\sffamily\fontsize{9.4}{11.5}\selectfont\color{Muted}#2}}}

\pagestyle{fancy}
\fancyhf{}
\fancyhead[L]{\sffamily\fontsize{9}{11}\selectfont\color{Muted}\CVName}
\fancyhead[R]{\sffamily\fontsize{9}{11}\selectfont\color{Muted}Curriculum vitae}
\fancyfoot[L]{\sffamily\fontsize{8.3}{10}\selectfont\color{Muted}Updated \CVUpdated}
\fancyfoot[R]{\sffamily\fontsize{8.3}{10}\selectfont\color{Muted}\thepage\ /\ \pageref*{LastPage}}
\renewcommand{\headrulewidth}{0pt}
\renewcommand{\footrulewidth}{0pt}
\fancypagestyle{firstpage}{\fancyhead{}}

\begin{document}
\thispagestyle{firstpage}

% ============================== HEADER ===================================
\begin{center}
  {\sffamily\bfseries\fontsize{27}{31}\selectfont\color{Accent}\CVName\par}
  \vspace{5pt}
  {\fontsize{10.2}{12.4}\selectfont
    \href{mailto:cli911@gatech.edu}{Email}\contactsep
    \href{https://lichangh20.github.io/}{Homepage}\contactsep
    \href{https://scholar.google.com/citations?hl=en\&user=9D8KEmYAAAAJ}{Google Scholar}\contactsep
    \href{https://github.com/lichangh20}{GitHub}\par}
\end{center}
\vspace{-3pt}

% ========================== RESEARCH INTERESTS ============================
\cvsection{Research Interests}
\textit{I work toward \textbf{Agentic Recursive Self-Improvement} to advance
\textbf{AutoResearch} and \textbf{Long-Horizon Problem Solving}.}\par
\vspace{3pt}
\smalltext{\textbf{Focus:} Harness Evolving; On-policy Learning;
Environments \& Rubrics; Agentic-training Infrastructure.}

% ============================= EDUCATION =================================
\cvsection{Education}
\cvrow{Georgia Institute of Technology}{Aug 2024 -- Present}
Ph.D. in Computational Science and Engineering\hfill GPA: 4.00/4.00\par
\smalltext{Advisors: \href{https://bo-dai.github.io/}{Bo Dai} and \href{https://chaozhang.org/}{Chao Zhang}}
\vspace{5pt}
\cvrow{Tsinghua University}{Sep 2020 -- Jun 2024}
B.Eng. in Computer Science and Technology\hfill GPA: 3.89/4.00\par
\smalltext{Advisors: \href{https://ml.cs.tsinghua.edu.cn/~jianfei/}{Jianfei Chen} and \href{https://ml.cs.tsinghua.edu.cn/~jun/index.shtml}{Jun Zhu}}

% ======================== INDUSTRY EXPERIENCE ============================
\cvsection{Industry Experience}
\job{ByteDance Seed}{Student Researcher}{Aug 2026 -- Present}
\smalltext{San Jose, CA\sep \textit{Worked with:} \href{https://simonxin.com/}{Xin Dong}}
\begin{cvitems}
  \item \textbf{Long-horizon coding agents:} Improving Seed models' behavior through post-training to strengthen performance on long-horizon coding tasks.
  \item \textbf{AutoResearch:} Building agents that identify worthwhile research questions and develop promising ideas, extending the research agenda from task execution to research judgment.
\end{cvitems}
% Draft: add concrete methods, ownership, and results only after author confirmation.
\vspace{5pt}
\job{Qwen}{Research Intern}{May 2026 -- Aug 2026}
\smalltext{Seattle, WA\sep \textit{Worked with:} \href{https://guoyin.wang/}{Guoyin Wang} and \href{https://fdias.fudan.edu.cn/fdiasen/ee/f4/c20151a192244/page.htm}{Jingren Zhou}}
\begin{cvitems}
  \item \textbf{Unified and efficient OPD:} Proposed a unified on-policy distillation framework organizing prior methods around \textbf{four design choices}, evaluated \textbf{30+ configurations}, and developed a simple, effective, and efficient OPD algorithm.
  \item \textbf{Large-scale OPD:} Contributed to the training infrastructure for \textbf{Qwen 3.6--3.8}, implementing OPD algorithms and validating their effectiveness on a \textbf{120B-parameter model} for long-horizon coding tasks, including SWE-bench and Terminal-Bench.
\end{cvitems}
% Research contributions: author-provided Revisit_OPD.pdf and Stochastic_On_policy_Distillation.pdf.
% Qwen 3.6--3.8 infrastructure, OPD implementation, and 120B validation confirmed by the author.
\vspace{5pt}
\job{AWS AI Lab}{Applied Scientist Intern}{May 2025 -- Aug 2025}
\smalltext{New York, NY\sep \textit{Worked with:} \href{https://mansimov.io/}{Elman Mansimov}, \href{https://www.amazon.science/author/elad-ben-avraham}{Elad Ben Avraham}, \href{https://www.amazon.science/author/ron-litman}{Ron Litman}, and \href{https://www.amazon.science/author/wei-xia-1}{Wei Xia}}
\begin{cvitems}
  \item \textbf{Deep Research evaluation:} Co-developed \href{https://aclanthology.org/2026.acl-long.448/}{\textbf{DREAM}}, a reference-free evaluation framework combining tool-based verification and adaptive rubrics to assess factual accuracy, temporal validity, and reasoning quality in Deep Research reports.
\end{cvitems}

% =========================== RESEARCH PROJECTS ============================
% Academic projects; industry research is described under the corresponding role.
% Personal contributions follow the author's explicit statements.
\cvsection{Selected Research Projects}
\project{2026}{Agent DAgger: on-policy learning for coding agents}{\href{https://arxiv.org/abs/2605.12913}{[Paper]}\enspace\href{https://github.com/STACX/stacx}{[Code]}}
\begin{projectitems}
  \item \textbf{Teacher-interleaved rollouts:} Combined on-policy state coverage with dense expert supervision to reduce covariate shift in long-horizon coding agents.
  \item \textbf{SWE-bench performance:} Achieved \textbf{27.3\% / 29.8\%} task resolution with 4B / 8B students on \mbox{SWE-bench Verified}, exceeding the strongest evaluated post-training baseline by \textbf{3.9 / 3.6 percentage points}.
\end{projectitems}

\project{2026}{STACX: infrastructure for end-to-end agentic RL}{\href{https://github.com/STACX/stacx}{[Code]}}
\begin{projectitems}
  \item \textbf{Project Lead:} Led overall agentic system planning, connecting sandboxed execution, rollout collection, verifier feedback, and model updates.
  \item \textbf{Modular agent training:} Supports interchangeable agents, environments, and learning algorithms within a shared training and evaluation workflow.
\end{projectitems}
% Project leadership and overall system planning confirmed by the author.

\project{2026}{Matryoshka Agent: learned orchestration for long-horizon MLE}{\href{https://arxiv.org/abs/2607.25090}{[Paper]}}
\begin{projectitems}
  \item \textbf{Hierarchical agent pipeline:} Helped connect high-level orchestration, sub-agent execution, and iterative ML experimentation for long-horizon tasks.
  \item \textbf{Small-model orchestration:} The framework enabled a \textbf{Qwen3-4B} orchestrator to approach o4-mini performance and improved Qwen3-30B-Coder by \textbf{up to 36.7\% relative} on the evaluated MLE tasks.
\end{projectitems}
% Personal contribution: helped build the long-horizon agentic pipeline.

\project{2025}{Matryoshka Pilot: learned guidance for black-box LLMs}{\href{https://arxiv.org/abs/2410.20749}{[Paper]}\enspace\href{https://github.com/lichangh20/Matryoshka}{[Code]}}
\begin{projectitems}
  \item \textbf{On-policy harness learning:} Developed a lightweight controller that learns multi-turn guidance from on-policy preference feedback to steer black-box LLMs.
  \item \textbf{Black-box adaptation:} Improved reasoning, planning, and personalization without updating the underlying model's weights.
\end{projectitems}

\project{2025}{MLE-Dojo: interactive environments for ML agents}{\href{https://arxiv.org/abs/2505.07782}{[Paper]}\enspace\href{https://github.com/MLE-Dojo/MLE-Dojo}{[Code]}}
\begin{projectitems}
  \item \textbf{Executable environments:} Helped build the environment framework and contributed task examples for iterative ML experimentation, debugging, and solution refinement.
  \item \textbf{Training and evaluation at scale:} The framework covers \textbf{200+ Kaggle competitions} with executable feedback; evaluations of \textbf{eight frontier LLMs} identified limitations in long-horizon problem solving.
\end{projectitems}
% Personal contributions: environment construction and selected task examples.

% =========================== PUBLICATIONS ================================
\cvsection{Publications \& Preprints}
\meta{\textsuperscript{*} Equal contribution.\enspace\textsuperscript{\textdagger} Equal second authorship.}

% Publications grouped by research theme, then year within each theme.
\pubtheme{Harness Evolution}

\publication
  {Matryoshka Agent: Unfolding Sub-Agents for Long-Horizon Machine Learning Engineering}
  {Rushi Qiang, \self\second, Haotian Sun\second, Yuchen Zhuang, Chao Zhang, Bo Dai.}
  {Preprint, 2026}
  {https://arxiv.org/abs/2607.25090}
  {}

\publication
  {Matryoshka Pilot: Learning to Drive Black-Box LLMs with LLMs}
  {\self\equal, Yuchen Zhuang\equal, Rushi Qiang, Haotian Sun, Hanjun Dai, Chao Zhang, Bo Dai.}
  {NeurIPS 2025}
  {https://arxiv.org/abs/2410.20749}
  {https://github.com/lichangh20/Matryoshka}

\pubtheme{Efficient On-Policy Learning}

\publication
  {Revisiting DAgger in the Era of LLM-Agents}
  {\self, Rushi Qiang\second, Jiawei Huang\second, Chenxiao Gao\second, Chao Zhang, Niao He, Bo Dai.}
  {Preprint, 2026}
  {https://arxiv.org/abs/2605.12913}
  {https://github.com/STACX/stacx}

\publication
  {Exploration-Driven Optimization for Test-Time Large Language Model Reasoning}
  {\self, Yuchen Zhuang, Chenxiao Gao, Haotian Sun, Rushi Qiang, Chao Zhang, Bo Dai.}
  {TMLR 2026}
  {https://arxiv.org/abs/2605.09853}
  {}

\pubtheme{Environments \& Rubrics}

\publication
  {DREAM: Deep Research Evaluation with Agentic Metrics}
  {Elad Ben Avraham\equal, \self\equal, Ron Dorfman\equal, Roy Ganz, Oren Nuriel, Amir Dudai, Aviad Aberdam, Noah Flynn, Elman Mansimov, Adi Kalyanpur, Ron Litman.}
  {ACL 2026 -- Oral}
  {https://aclanthology.org/2026.acl-long.448/}
  {}

\publication
  {MLE-Dojo: Interactive Environments for Empowering LLM Agents in \mbox{Machine Learning Engineering}}
  {Rushi Qiang\equal, Yuchen Zhuang\equal, Yinghao Li, Dingu Sagar V K, Rongzhi Zhang, \self, Ian Shu-Hei Wong, Sherry Yang, Percy Liang, Chao Zhang, Bo Dai.}
  {NeurIPS 2025}
  {https://arxiv.org/abs/2505.07782}
  {https://github.com/MLE-Dojo/MLE-Dojo}

\pubtheme{Model Efficiency}

\publication
  {Task-Specific Zero-shot Quantization-Aware Training for Object Detection}
  {\self\equal, Xinrui Chen\equal, Ji Wang\equal, Kang Zhao, Jianfei Chen.}
  {ICCV 2025}
  {https://arxiv.org/abs/2507.16782}
  {https://github.com/DFQ-Dojo/dfq-toolkit}

\publication
  {Training Transformers with 4-bit Integers}
  {Haocheng Xi, \self, Jianfei Chen, Jun Zhu.}
  {NeurIPS 2023}
  {https://arxiv.org/abs/2306.11987}
  {https://github.com/haochengxi/Train_Transformers_with_INT4}

% ============================ AWARDS / SERVICE ===========================
\cvsection{Selected Honors}
\award{\textbf{Outstanding Graduate Award}, Tsinghua University}{2024}
\award{\textbf{Comprehensive Excellence Scholarship}, Tsinghua University}{2023}
\award{\textbf{Academic \& Social Service Excellence Scholarship}, Tsinghua University}{2022}
% The 2020 high-school award is preserved in EDITORIAL_NOTES.md, omitted here
% to prioritize more recent achievements in the two-page industry version.

\cvsection{Academic Service}
\smalltext{\textbf{Reviewer:} NeurIPS, ICLR, ICML, ACL Rolling Review, TMLR.}
\smalltext{\textbf{Teaching Assistant:} \href{https://bo-dai.github.io/CX4240-spring2026/}{CX4240 -- Computing for Data Analysis}, Georgia Tech\hfill{\sffamily\fontsize{9.4}{11.5}\selectfont\color{Muted}Spring 2026}}
\cvsection{Technical Skills}
\smalltext{\textbf{Languages \& frameworks:} Python, C++, PyTorch, CUDA.}
\smalltext{\textbf{Systems experience:} Implemented CUDA kernels for 4-bit transformer training (NeurIPS 2023).}

\end{document}
